\section{Larger regularity neighborhoods of the amplitude segment}
\label{ext:section}

The convergence theorem above and the extension in this section have different
conclusions. Here the second corrected reference gives larger sufficient
initial-data classes for global regularity and a uniform Galerkin budget.
No convergence of the ordinary heat-seeded Picard iteration is asserted on
these larger classes.

In this section $F,G,C,H$ denote the fixed spatial Fourier fields and
$L_1,L_5,L_6$ their scalar heat profiles. The initial radius is $r$, the
critical-error barrier is $B<2$, and $P$ is the reference action. We write
$M_{\mathrm{gate}}$ for the scalar endpoint test and $\Gamma$ for the total
solution-action bound; the integer $L$ is its rounded budget exponent.
The scalar barrier $B$ is distinct from the bilinear operators
$B(\cdot,\cdot)$ and $B_s(\cdot,\cdot)$.

\begin{theorem}[Regularity regions with explicit budgets]\label{ext:theorem}
Let $u_0=aW+h$, where $a\in\mathbb R$ and $h$ is smooth, real, mean-zero,
and divergence-free. Suppose that one row of Table~\ref{ext:regions} satisfies
$|a|\le a_*$ and $A_{-1}(h)\le r$. Then the viscosity-one unforced
Navier--Stokes equation has a unique global smooth solution with datum $u_0$.
Writing
\[
 E_0=\|u_0\|_2^2,\qquad X_0=\|\nabla u_0\|_2^2,
\]
its spherical Galerkin solutions satisfy
\begin{align}
 \sup_{N,t\ge0}\|\nabla u_N(t)\|_2^2&\le X_0e^L,
 \label{ext:h1}\\
 \sup_N\int_0^\infty\|\Delta u_N(t)\|_2^2\,dt&\le X_0e^L,
 \label{ext:h2}\\
 B_*(u_0):=\sup_N\int_0^\infty\|\nabla u_N(t)\|_2^4\,dt
 &\le\frac{E_0X_0}{2}e^L.
 \label{ext:nsb}
\end{align}
There is no restriction on the Fourier support of $h$.
\end{theorem}

\begin{table}[htbp]
\centering
\begin{tabular}{ccccc}
\hline
Amplitude $a_*$ & Radius $r$ & Error barrier $B$ & $\Gamma$ (orientation) & $L$\\
\hline
$1$ & $2/25$ & $6/5$ & $24.40770\ldots$ & $25$\\
$1$ & $1/10$ & $9/5$ & $58.60140\ldots$ & $59$\\
$11/10$ & $1/40$ & $6/5$ & $28.06654\ldots$ & $29$\\
$9/8$ & $1/100$ & $6/5$ & $29.06160\ldots$ & $30$\\
$8/7$ & $1/100$ & $8/5$ & $48.27248\ldots$ & $49$\\
\hline
\end{tabular}
\caption{Five certified amplitude--radius--budget triples. The first and
fourth rows are nested refinements of the second and fifth, respectively,
with smaller budget exponents. The displayed values of $\Gamma$, defined
in \eqref{ext:G}, are for orientation; the exact rational enclosures
prove $\Gamma<L$.}
\label{ext:regions}
\end{table}

Write $\mathcal C(A,r)$ for the set of the stated smooth data admitting a
decomposition $u_0=aW+h$ with $|a|\le A$ and $A_{-1}(h)\le r$. Then
\[
 \mathcal C(1,2/25)\subset\mathcal C(1,1/10),\qquad
 \mathcal C(9/8,1/100)\subset\mathcal C(8/7,1/100).
\]
Thus the five rows have three inclusion-maximal data classes among the
listed classes. The two nested rows retain their value through the
smaller exponents $25$ instead of $59$ and $30$ instead of $49$.
For a datum belonging to more than one row, one may use the smallest
listed exponent for which an admissible decomposition exists. These
observations do not certify a new amplitude--radius rectangle formed
by mixing bounds from different rows, and no maximality among all
regularity classes is asserted. The quantities $E_0$ and $X_0$ always
refer to the full datum; the critical radius alone does not uniformly
bound its $H^1$ norm.

\subsection{The complete critical error inequality}
\label{ext:comparison}

Let $v$ be a smooth finite Fourier reference and set
\[
 r_v=\partial_tv-\Delta v+B_s(v,v),\qquad
 U(t)=A_0(v(t)),\qquad R(t)\ge A_{-1}(r_v(t)).
\]
At every cutoff containing the support of $v$, the exact Galerkin error
$e=u_N-v$ satisfies
\[
 \partial_te-\Delta e+
 \Pi_N\{2B_s(v,e)+B_s(e,e)\}=-\Pi_Nr_v.
\]
Here and below $\Pi_N$ is the spherical projection. The residual is the full
Fourier residual before projection; no exterior modes have been discarded
when defining $R$.

Put $b=A_{-1}(e)$ and $d=A_1(e)$. The symmetric Leray estimate and
$A_0(e)^2\le bd$ give, almost everywhere,
\begin{equation}
 b'+(1-b/2)d\le U\sqrt{bd}+R.
 \label{ext:error}
\end{equation}
Indeed, the two mixed terms contribute at most $UA_0(e)$ and the self term
at most $A_0(e)^2/2$. The finite coefficient norms are absolutely continuous;
the same inequalities hold at coefficient zeros by the integrated norm
inequality, or by regularization and passage to the limit.

\begin{lemma}[Endpoint gate and retained history]\label{ext:historylemma}
Suppose
\[
 b(0)\le r<B<2,\qquad
 P=\int_0^\infty U(t)^2\,dt<\infty,\qquad
 R_*=\int_0^\infty R(t)\,dt<\infty,
\]
where $R\ge0$. Define $\Psi(s)=se^{-s/2}$ for $0\le s<2$.
If
\begin{equation}
 M_{\mathrm{gate}}:=e^{P/4}\Psi(r)+
 \int_0^\infty
  \exp\!\left(\frac14\int_t^\infty U(s)^2\,ds\right)R(t)\,dt
 <\Psi(B),
 \label{ext:gate}
\end{equation}
then $b(t)<B$ for every $t\ge0$. With
\begin{equation}
 F_*(s)=-2s-4\log(1-s/2),\qquad
 C_*=F_*(r)+\frac{B}{1-B/2}R_*,
 \label{ext:F}
\end{equation}
one moreover has
\begin{align}
 \left(\int_0^\infty b(t)d(t)\,dt\right)^{1/2}
 &\le\sqrt{F_*(B)}+\sqrt{F_*(B)+C_*},
 \label{ext:errorbudget}\\
 \int_0^\infty A_0(u_N(t))^2\,dt
 &\le
 \left[\sqrt P+\sqrt{F_*(B)}+\sqrt{F_*(B)+C_*}\right]^2.
 \label{ext:referencebudget}
\end{align}
\end{lemma}

\begin{proof}
While $b<2$, maximize $U\sqrt b\sqrt d-(1-b/2)d$ in $\sqrt d\ge0$.
As $0<\Psi'(b)=e^{-b/2}(1-b/2)\le1$, this gives
\begin{equation}
 (\Psi(b))'\le\frac14U^2\Psi(b)+R.
 \label{ext:transformed}
\end{equation}
Let $J(t)=\int_0^tU^2$. Integrating \eqref{ext:transformed} only up to $t$
and multiplying by $e^{(P-J(t))/4}$ yields
\[
 e^{(P-J(t))/4}\Psi(b(t))
 \le e^{P/4}\Psi(r)+
 \int_0^t e^{(P-J(s))/4}R(s)\,ds\le M_{\mathrm{gate}}.
\]
If $b$ first reached $B$, the left side at that time would be at least
$\Psi(B)$, contradicting \eqref{ext:gate}. Finite Galerkin solutions exist
globally by their kinetic-energy identity, so this first-contact argument
proves $b<B$ on the entire half-line without assuming a global PDE solution.
It also proves the stronger envelope
\begin{equation}
 \Psi(b(t))\le\Psi(B)e^{-(P-J(t))/4}.
 \label{ext:envelope}
\end{equation}

Define $\beta(t)\in[0,B]$ by equality in \eqref{ext:envelope}, using only
the increasing branch of $\Psi$ on $[0,2)$. Then $b\le\beta\le B$ and
\[
 \beta'=\frac{U^2\beta}{4(1-\beta/2)},\qquad
 F_*'(s)=\frac{s}{1-s/2}.
\]
The chain rule therefore proves, for every finite $T$,
\begin{equation}
 \int_0^T U^2\frac{b^2}{(1-b/2)^2}\,dt
 \le4\{F_*(\beta(T))-F_*(\beta(0))\}\le4F_*(B).
 \label{ext:weightedhistory}
\end{equation}
No strict monotonicity of $J$ is required. In particular intervals on which
$U=0$ cause no change-of-variable problem.

Multiply \eqref{ext:error} by $b/(1-b/2)$. After integration on $[0,T]$,
drop the nonnegative terminal value $F_*(b(T))$. For
$Q_T=\int_0^Tbd$, monotonicity of $s/(1-s/2)$ and
\eqref{ext:weightedhistory} give
\[
 Q_T\le F_*(r)+\frac{B}{1-B/2}R_*+
 \left(4F_*(B)Q_T\right)^{1/2}.
\]
The positive root of this quadratic inequality bounds $\sqrt{Q_T}$ by
the right side of \eqref{ext:errorbudget}, uniformly in $T$.
Only then take $T\to\infty$ by monotone convergence. Finally,
$A_0(e)^2\le bd$ and Minkowski in time prove
\eqref{ext:referencebudget}.
\end{proof}

\paragraph{Reference-free baseline.}
For smooth, real, mean-zero, solenoidal data, put $r_0=A_{-1}(u_0)<2$,
choose $r_0<B<2$, and take $v\equiv0$. Then $U=R=P=0$, and the gate is
$\Psi(r_0)<\Psi(B)$, which holds because $\Psi$ is strictly increasing
on $[0,2)$. The zero reference belongs to every Galerkin subspace, and
$A_{-1}(\Pi_Nu_0)\le r_0$. Lemma~\ref{ext:historylemma}, with $r=r_0$,
therefore gives a finite action bound uniform in every cutoff.
The enstrophy estimate and the compactness, uniqueness and continuation
argument of Section~\ref{ext:pde} then give a global smooth solution.
This is a consequence of the present comparison estimate; the threshold
$2$ is not attributed to the earlier small-data theorems.

\subsection{The second corrected reference and its signed residual}
\label{ext:reference}

All fields in this subsection are finite Fourier polynomials. Set
\begin{align*}
 F&=B_s(W,W),&
 G&=2B_s(W,F)=-\tfrac12W+C,\\
 C&=\tfrac12(\sin z\cos2y,\sin x\cos2z,\sin y\cos2x),&
 H&=B_s(F,F).
\end{align*}
The frequency shells of $W,F,C,H$ have squared radii $1,2,5,6$,
respectively. Define
\begin{align}
 L_1(t)&=\frac{e^{-t}-(1+2t)e^{-3t}}4,&
 L_5(t)&=\frac{e^{-5t}+(2t-1)e^{-3t}}4,\notag\\
 L_6(t)&=\frac{(2t^2-2t+1)e^{-4t}-e^{-6t}}4.
 \label{ext:profiles}
\end{align}
Their initial values vanish and
\[
 (\partial_t+1)L_1=(\partial_t+5)L_5=te^{-3t},\qquad
 (\partial_t+6)L_6=t^2e^{-4t}.
\]
Thus they are nonnegative causal heat convolutions. The reference is
\begin{equation}
 w_{2,a}=(ae^{-t}-a^3L_1/2)W-a^2te^{-2t}F
                +a^3L_5C-a^4L_6H.
 \label{ext:w2}
\end{equation}
It starts exactly at $aW$ and has support $|k|\le\sqrt6$.
For clarity, if $w_1=ae^{-t}W-a^2te^{-2t}F$ and
$z_2=(a^3/2)L_1W-a^3L_5C+a^4L_6H$, then $w_{2,a}=w_1-z_2$,
\[
 r_1=-a^3te^{-3t}G+a^4t^2e^{-4t}H,
 \qquad (\partial_t-\Delta)z_2=r_1.
\]
It follows algebraically that the full residual of \eqref{ext:w2} is
$-2B_s(w_1,z_2)+B_s(z_2,z_2)$.

For $m=|a|$, introduce the nonnegative scalar functions
\[
 x=me^{-t},\quad y=m^2te^{-2t},\quad
 p=m^3L_1/2,\quad q=m^3L_5,\quad s=m^4L_6,
\]
and put $K_{AB}=2B_s(A,B)$. For $a=m\ge0$, the complete residual is
\begin{align}
 r_2={}&(p^2-2xp)F+ypG+(x-p)qK_{WC}-(x-p)sK_{WH}\notag\\
       &-yqK_{FC}+ysK_{FH}+q^2B_s(C,C)
                       -qsK_{CH}+s^2B_s(H,H).
 \label{ext:residual}
\end{align}
For $a<0$, the corresponding signed coefficients have exactly the same
magnitudes as those used below. This follows by substituting the powers
$a,a^2,a^3,a^4$ before grouping. It does not assume that the full solution
is invariant under $u\mapsto-u$.

The finite Fourier identities and upper bounds required here are recorded
in Table~\ref{ext:caps}. Each bound is for the complete generated Fourier
field with the Leray projection, not for its projection back to the original
support. Their coefficient-level verification is part of the certificate
appendix.

\begin{table}[htbp]
\centering
\begin{tabular}{lrr}
\hline
Field & $A_{-1}$ upper cap $c_j$ & Nonzero modes\\
\hline
$F$ & $17/8$ & $12$\\
$K_{WC}$ & $171/100$ & $48$\\
$K_{WH}$ & $27/20$ & $72$\\
$G$ & $87/40$ & $18$\\
$K_{FC}$ & $6/5$ & $60$\\
$K_{FH}$ & $183/100$ & $120$\\
$B_s(C,C)$ & $11/25$ & $48$\\
$K_{CH}$ & $3/4$ & $162$\\
$B_s(H,H)$ & $9/25$ & $120$\\
\hline
\end{tabular}
\caption{Complete Fourier norm caps in the order used in \eqref{ext:R}.}
\label{ext:caps}
\end{table}

Since $p/x\le m^2/8$, all entries in the following vector are nonnegative
for $m^2\le3$:
\begin{equation}
 (f_1,\ldots,f_9)=
 (2xp-p^2,(x-p)q,(x-p)s,yp,yq,ys,q^2,qs,s^2),
 \qquad R_m=\sum_{j=1}^9c_jf_j.
 \label{ext:R}
\end{equation}
Consequently $A_{-1}(r_2)\le R_m$. The subtractive terms in the first three
profiles remain grouped before a Fourier norm bound is taken.

\subsection{Why the endpoint calculations cover amplitude intervals}
\label{ext:monotonicity}

The four shells in \eqref{ext:w2} are disjoint, and their $A_0$ norms are
$3,3,3/2,\sqrt3$. For $m^2\le8$, the fundamental coefficient retains the
sign of $a$, since $L_1\le e^{-t}/4$. Hence the exact reference norm is
\begin{equation}
 U_2(m,t)=3me^{-t}+3m^2te^{-2t}
      -\frac32m^3(L_1-L_5)+\sqrt3m^4L_6.
 \label{ext:U2}
\end{equation}
For $k_\lambda(t)=e^{-\lambda t}\mathbf1_{t\ge0}$, direct convolution
identities give
\[
 L_1-L_5=4k_1*k_3*k_3*k_5,\qquad
 2L_6=4k_4*k_4*k_4*k_6.
\]
Comparison of the four nonnegative factors gives $L_1-L_5\ge2L_6$.
Together with $0\le L_1-L_5\le L_1\le e^{-t}/4$, this proves
\begin{align}
 U_2(m,t)&\le U_1(m,t):=3me^{-t}+3m^2te^{-2t}
                   &&(m^2\le3),\label{ext:U1}\\
 \partial_mU_2(m,t)&\ge(3-9m^2/8)e^{-t}\ge0
                   &&(m^2\le8/3).
 \label{ext:Umonotone}
\end{align}
In the second line the omitted terms are nonnegative. Replacing $\sqrt3$
by a rational upper bound preserves this monotonicity, as $L_6\ge0$.
Thus a reference-budget upper enclosure computed at $a_*$ bounds the
budget throughout $0\le m\le a_*$.

The residual majorant is monotone on the same required interval.
The first three profiles in \eqref{ext:R} have the forms
$m^4A-m^6D$, $m^4A-m^6D$, and $m^5A-m^7D$, respectively, with
$D/A\le1/16,1/8,1/8$ wherever $A>0$. Their derivatives have nonnegative
factors bounded below by
\[
 4-6m^2/16,\qquad4-6m^2/8,\qquad5-7m^2/8,
\]
all strictly positive for $m^2\le3$. Where $A=0$, the corresponding
profile vanishes and the same conclusion follows by continuity. The six
remaining profiles are nonnegative pure powers of $m$. Every amplitude
endpoint in Table~\ref{ext:regions} lies in both stated ranges.

\subsection{Finite exact inequalities behind the five gates}
\label{ext:arithmetic}

Here we specify the finite arithmetic used in \eqref{ext:gate}, so that
the application is not an invocation of a bare program status. A finite
exponential polynomial $T(t)=\sum_{\lambda,d}c_{\lambda d}t^de^{-\lambda t}$
with rational coefficients is integrated by
\begin{equation}
 \mathcal I_n(T):=\int_0^\infty e^{-2nt}T(t)\,dt
 =\sum_{\lambda,d}c_{\lambda d}
             \frac{d!}{(\lambda+2n)^{d+1}}.
 \label{ext:moments}
\end{equation}
Use the explicit upper enclosure
\[
 \sqrt3<s_3^+:=\frac{866025403784438647}{500000000000000000},
\]
whose square exceeds $3$. Let $U_2^+(a_*,t)$ denote \eqref{ext:U2} with
$\sqrt3$ replaced by $s_3^+$. It is nonnegative and dominates $U_2(m,t)$
for $0\le m\le a_*$. Set
\begin{equation}
 P^+(a_*)=\mathcal I_0((U_2^+(a_*,\cdot))^2),\qquad
 R_*(a_*)=\mathcal I_0(R_{a_*}),\qquad
 I_n(a_*)=\mathcal I_n(R_{a_*}).
 \label{ext:certinputs}
\end{equation}
These are exact rational numbers. Although the expanded polynomial $R_m$
has some negative coefficients, \eqref{ext:R} proves $R_m(t)\ge0$;
therefore $0\le I_{n+1}(m)\le I_n(m)$.

By \eqref{ext:U1},
\[
 \int_0^\infty U_1(m,t)^2\,dt
 =\frac92m^2+2m^3+\frac9{32}m^4.
\]
Subtracting only the heat part of the \emph{past} $U_1^2$ integral gives
\begin{equation}
 \exp\!\left(\frac14\int_t^\infty U_2(m,s)^2\,ds\right)
 \le e^{\Delta_m}\exp(c_me^{-2t}),\qquad
 c_m=\frac98m^2,\quad \Delta_m=\frac12m^3+\frac9{128}m^4.
 \label{ext:weight}
\end{equation}
This uses $U_1\ge3me^{-t}$, not a lower bound of that form for $U_2$.
Tonelli's theorem applies to the positive exponential series against $R_m$.
If $t_n=c_m^nI_n(m)/n!$, then
$t_{n+1}\le c_mt_n/(n+1)$, so
\begin{equation}
 \int_0^\infty e^{\frac14\int_t^\infty U_2^2}R_m(t)\,dt
 \le e^{\Delta_m}
 \left(\sum_{n=0}^{24}t_n+\frac{t_{25}}{1-c_m/26}\right).
 \label{ext:tail}
\end{equation}
All denominators are positive at the stated endpoints.

For rational $0\le z<66$ define the rational upper bound
\begin{equation}
 \mathcal E(z)=\sum_{n=0}^{64}\frac{z^n}{n!}
       +\frac{z^{65}/65!}{1-z/66}\ge e^z.
 \label{ext:expupper}
\end{equation}
For a given row, replace the initial term and residual term in the gate by
\begin{align}
 M_{\mathrm{gate}}^+(a_*,r)={}&r\mathcal E(P^+(a_*)/4-r/2)\notag\\
 &+\mathcal E(\Delta_{a_*})
 \left(\sum_{n=0}^{24}\frac{c_{a_*}^n}{n!}I_n(a_*)+
       \frac{c_{a_*}^{25}I_{25}(a_*)/25!}{1-c_{a_*}/26}\right),
 \label{ext:Mplus}\\
 \Psi^-(B)={}&\frac{B}{\mathcal E(B/2)}\le\Psi(B).
 \label{ext:psilower}
\end{align}
The argument $P^+/4-r/2$ is nonnegative in every row. Monotonicity proved
above, and monotonicity of $\Psi$ on the initial interval, imply
$M_{\mathrm{gate}}\le M_{\mathrm{gate}}^+$
for every datum in that row. Direct rational evaluation of
\eqref{ext:moments}--\eqref{ext:psilower} gives the strict inequalities
recorded in Table~\ref{ext:gatetable}. The endpoints of the data classes
are closed because the gate inequalities themselves have strict margins.

\begin{table}[htbp]
\centering
\begin{tabular}{ccccc}
\hline
$a_*$ & $r$ & $B$ & $M_{\mathrm{gate}}^+$ (orientation) &
 $\Psi^-(B)-M_{\mathrm{gate}}^+$ (orientation)\\
\hline
$1$ & $2/25$ & $6/5$ & $0.636555727041\ldots$ & $0.02201823627\ldots$\\
$1$ & $1/10$ & $9/5$ & $0.730918119157\ldots$ & $0.00090726837\ldots$\\
$11/10$ & $1/40$ & $6/5$ & $0.656895071385\ldots$ & $0.00167889192\ldots$\\
$9/8$ & $1/100$ & $6/5$ & $0.636557679768\ldots$ & $0.02201628354\ldots$\\
$8/7$ & $1/100$ & $8/5$ & $0.714761998246\ldots$ & $0.00416434434\ldots$\\
\hline
\end{tabular}
\caption{The exact numbers are the rational expressions
\eqref{ext:Mplus}--\eqref{ext:psilower}, stored in the endpoint certificate.
The displayed decimals are explanatory; the finite integer comparisons
prove that each margin is positive.}
\label{ext:gatetable}
\end{table}

For each passing row define
\begin{equation}
 \Gamma=\left[\sqrt{P^+(a_*)}+\sqrt{F_*(B)}+
 \sqrt{F_*(B)+F_*(r)+\frac{B}{1-B/2}R_*(a_*)}\right]^2.
 \label{ext:G}
\end{equation}
Its strict upper comparison $\Gamma<L$ in Table~\ref{ext:regions} is another
finite certificate, rather than a floating-point evaluation of logarithms.
Explicitly, normalize each logarithm argument as $2^kx$ with $1\le x\le2$.
For $q=(x-1)/(x+1)\in[0,1/3]$,
\[
 2\sum_{j=0}^{39}\frac{q^{2j+1}}{2j+1}
 \le\log x\le
 2\sum_{j=0}^{39}\frac{q^{2j+1}}{2j+1}
       +\frac{2q^{81}}{81(1-q^2)}.
\]
Apply the same enclosure to $\log2$. Square roots of the resulting
nonnegative rationals are bounded above at scale $10^{30}$ by integer
square roots with outward rounding. Substitution in \eqref{ext:G}
then compares rational numbers with the listed integers $L$.

The exact certificate dependencies for this section are:
\begin{center}
\begin{tabular}{p{0.41\textwidth}p{0.51\textwidth}}
\hline
Certificate & Mathematical role\\
\hline
\path{second_correction_fourier_check.py}
 & Fourier fields, profiles, projected products, and the nine caps in
 Table~\ref{ext:caps}.\\
\path{enlarged_endpoint_gate.py}
 & Rational moments, $P^+$, the full weighted residual, and the five
 strict endpoint inequalities.\\
\path{enlarged_budget_static.py}
 & Outward logarithm and square-root enclosures proving the five
 inequalities $\Gamma<L$.\\
\hline
\end{tabular}
\end{center}
The accompanying exact output files and hashes are supplied in the
certificate appendix. The analytic implications of these finite
inequalities are established here; a hash alone is not a mathematical
verification of a norm estimate or a PDE passage to the limit.

\subsection{All cutoffs and the full PDE conclusion}
\label{ext:pde}

\begin{proof}[Proof of Theorem~\ref{ext:theorem}]
Fix one row and its datum. At every cutoff $N\ge\sqrt6$, the reference
\eqref{ext:w2} is retained exactly. Its initial error is $\Pi_Nh$, and
projection contracts both $A_{-1}(h)$ and the full residual norm. The
preceding monotonicity and rational inequalities therefore verify the
hypotheses of Lemma~\ref{ext:historylemma} uniformly in $N$. In particular,
\begin{equation}
 \int_0^\infty A_0(u_N(t))^2\,dt\le\Gamma<L
 \qquad(N\ge\sqrt6).
 \label{ext:largecutaction}
\end{equation}

Write $X_N=\|\nabla u_N\|_2^2$ and $Y_N=\|\Delta u_N\|_2^2$.
Testing the exact Galerkin equation against $-\Delta u_N$, using
$\|u_N\|_\infty\le A_0(u_N)$, and applying $2ab\le a^2+b^2$ gives
\begin{equation}
 X_N'+Y_N\le A_0(u_N)^2X_N.
 \label{ext:enstrophy}
\end{equation}
The integrating factor for this scalar inequality and
\eqref{ext:largecutaction} prove
\[
 \sup_tX_N(t)\le X_0e^{\Gamma},\qquad
 \int_0^\infty Y_N(t)\,dt\le X_0e^{\Gamma}.
\]
The kinetic-energy equality gives $\int_0^\infty X_N\,dt\le E_0/2$;
hence $\int_0^\infty X_N^2\,dt\le(E_0X_0/2)e^{\Gamma}$.

For the lower cutoffs $N<\sqrt6$, no comparison with a partially removed
reference is needed. Directly from energy and the spectral support,
\begin{align*}
 \sup_tX_N(t)&\le N^2E_0\le6X_0,\\
 \int_0^\infty Y_N(t)\,dt
   &\le N^2\int_0^\infty X_N(t)\,dt\le3E_0\le3X_0,\\
 \int_0^\infty X_N(t)^2\,dt
   &\le\frac{N^2E_0^2}{2}\le3E_0X_0.
\end{align*}
Here $E_0\le X_0$ is the mean-zero Poincar\'e inequality; it does not
require a small $E_0$. Every listed $L\ge2$, and
$e^2>1+2+2+8/6=19/3>6$. These lower-cutoff bounds and $\Gamma<L$ prove
\eqref{ext:h1}--\eqref{ext:nsb} for every cutoff. We have not claimed
\eqref{ext:largecutaction} for cutoffs missing reference modes.

For completeness, the finite systems themselves exist for every time:
their energy identity bounds their finite coefficient vectors and prevents
finite-time escape in the polynomial ODE. On each finite time interval,
\eqref{ext:h1}--\eqref{ext:h2} bound the Galerkin sequence in
$L^\infty H^1\cap L^2H^2$. The equation bounds its time derivatives in
$L^2L^2$: by the periodic embedding $H^2\hookrightarrow L^\infty$,
\[
 \|(u_N\cdot\nabla)u_N\|_2
 \le\|u_N\|_\infty\|\nabla u_N\|_2
 \le C\|u_N\|_{H^2}\|u_N\|_{H^1}.
\]
Simon~\cite{Simon}, Corollary~4, therefore gives a subsequence strongly convergent in $L^2H^1$
on each finite interval, with the nonlinear term converging in
 distributions. The limit has the prescribed datum and is a global strong
solution. The bounds pass by weak lower semicontinuity. Applying Fatou
first to finite Fourier sums and then to their increasing limit also
preserves $\int_0^\infty A_0(u)^2\,dt\le\Gamma$ for the PDE solution.

To see uniqueness in this class, take the difference of two strong
solutions and test its equation in $L^2$. Incompressibility eliminates
one transport term, and integration by parts bounds the other by
$\|u\|_\infty\|w\|_2\|\nabla w\|_2$ for one of the solutions $u$.
Young's inequality and the locally integrable coefficient
$\|u\|_\infty^2$ give uniqueness by Gronwall. The local smooth solution
with smooth datum agrees with this strong solution. The uniform $H^1$
continuation bound and the local theory of \cite[Proposition~2.2]{Tao2009} extend smoothness
through every finite time. No global full-PDE regularity assumption was
used to obtain the Galerkin estimates.
\end{proof}

\paragraph{Scope of the extension.}
These are sufficient regions with explicit constants, not maximal regions
or an arbitrary-data assertion. They extend the initial-data classes covered
by the particular reference calculation while retaining a separate proof
from the ordinary-Picard convergence theorem. A failed upper gate would
show only that this comparison was inconclusive. Conditional criteria,
fieldwise modified-energy obstructions, and numerical experiments in other
branches do not enlarge Theorem~\ref{ext:theorem} without their own
unconditional full-PDE argument.
